\section{The equation the surface must obey}
\label{sec:lveq}

\subsection{Forward: from the field to every marginal}

Fix a bounded measurable field $v\ge0$ and let $Y$ solve \eqref{eq:lvsde}.
Everything follows from watching the call price evolve in maturity.

\begin{proposition}[The forward Dupire equation]
\label{prop:lvforward}
Suppose $Y_\tau$ has a density $f(\tau,\cdot)$ such that
$v y^{2}f$ and $y\,\partial_y(v y^{2}f)$ are integrable and vanish at
$y=0$ and as $y\to\infty$.
Then the normalized call $c(\tau,y)=\E[(Y_\tau-y)^+]$ satisfies
\begin{equation}
\boxed{\;\partial_\tau c(\tau,y)
=\tfrac12\,v(\tau,y)\,y^{2}\,\partial_{yy}c(\tau,y),
\qquad c(0,y)=(1-y)^{+}.\;}
\label{eq:lvdupire}
\end{equation}
\end{proposition}

\begin{proof}
By It\^o's formula, the density satisfies the Fokker--Planck equation
$\partial_\tau f=\tfrac12\,\partial_{yy}\!\big(v y^{2}f\big)$ in the weak
sense.  Integrate it against the payoff:
\[
\partial_\tau c(\tau,y_0)
=\int_0^\infty (y-y_0)^{+}\,\partial_\tau f(\tau,y)\,\dd y
=\tfrac12\int_0^\infty (y-y_0)^{+}\,
  \partial_{yy}\big(v y^{2}f\big)\,\dd y .
\]
Integrating by parts twice moves the two derivatives onto the payoff,
whose derivatives in $y$ are the step $\1_{\{y>y_0\}}$ and then the point
mass $\delta_{y_0}$, while the boundary terms vanish by the decay
hypothesis.  What remains is
$\tfrac12\,v(\tau,y_0)\,y_0^{2}f(\tau,y_0)$, and
$f(\tau,y_0)=\partial_{yy}c(\tau,y_0)$ is the density identity of
\cref{sec:intro}.  The initial condition is $Y_0=1$.
\end{proof}

Equation \eqref{eq:lvdupire} is a \emph{forward} parabolic equation in the
strike variable with the payoff as initial data: heat flow, with
$\tfrac12 v y^{2}$ as a spatially varying conductivity.  One field, one
linear march, and every marginal law of the surface comes out at once ---
this is the computational engine of the whole chapter.  The analogy with
heat is not decoration; it is the source of every structural property we
will use.  Heat flow spreads bumps and never sharpens them, and spreading,
translated into option language, is exactly the no-arbitrage direction:

\begin{proposition}[Positive local variance propagates no-arbitrage]
\label{prop:lvpropagate}
Let $c$ solve \eqref{eq:lvdupire} with $v\ge0$ bounded.  Then for every
$\tau>0$: (i) $\partial_{yy}c(\tau,\cdot)\ge0$ --- each slice is
butterfly-free; (ii) $\partial_\tau c(\tau,y)\ge0$ --- the normalized call
is non-decreasing in maturity at fixed $y$, which is the calendar order of
\cref{sec:calendar}.
\end{proposition}

\begin{proof}[Proof sketch]
Differentiating \eqref{eq:lvdupire} twice in $y$ shows that
$f=\partial_{yy}c$ satisfies
$\partial_\tau f=\tfrac12\,\partial_{yy}(v y^{2}f)$ --- the
Fokker--Planck equation again --- with initial datum
$\partial_{yy}(1-y)^{+}=\delta_{1}\ge0$.  Such equations preserve
nonnegativity: probabilistically, $f(\tau,\cdot)$ \emph{is} the density of
the diffusion \eqref{eq:lvsde} started at $1$; analytically, apply the
parabolic maximum principle to the adjoint equation.  With $f\ge0$
established, (ii) is \eqref{eq:lvdupire} itself:
$\partial_\tau c=\tfrac12 v y^{2}f\ge0$.
\end{proof}

Pause on what this proposition buys, because it is the entire strategy of
the chapter.  In Chapters~2 and~3, validity was work: the LQD chart made
butterfly freedom an identity within one slice, the volatility chart had to
police $g_{\rm D}\ge0$, and calendar order across slices was in both cases
a separate condition.  Here, positivity of \emph{one} function of two
variables propagates \emph{both} families across the whole surface, for
free, forever.  The rest of the chapter is the discipline of not
squandering that theorem: choose a parametrization in which $v>0$ is
trivial to enforce (\cref{sec:lvsheet}), choose a discretization that
inherits the maximum principle rather than merely approximating it
(\cref{sec:lvlattice}), and then measure what little the finite grid leaks
instead of assuming the theorem carried over.

\subsection{Backward: from the surface to the field}

Read \eqref{eq:lvdupire} the other way.  If the surface is given, the field
is determined pointwise,
\begin{equation}
v(\tau,y)=\frac{\partial_\tau c(\tau,y)}
{\tfrac12\,y^{2}\,\partial_{yy}c(\tau,y)},
\label{eq:lvextract}
\end{equation}
wherever the density in the denominator is positive.  This is the promised
uniqueness: within the class \eqref{eq:lvsde}, at most one field is
consistent with a given surface, and \eqref{eq:lvextract} computes it.  The
same formula becomes most instructive in the volatility chart.

\begin{proposition}[The field in the volatility chart]
\label{prop:lvratio}
Write the surface in implied form, $c=\Black(k,w(k,\tau))$ with
$w(\cdot,\tau)$ the total-variance slice.  Then \eqref{eq:lvextract} reads
\begin{equation}
v(\tau,e^{k})=\frac{\partial_\tau w(k,\tau)}{g_{\rm D}(k)},
\label{eq:lvratio}
\end{equation}
with $g_{\rm D}$ the Durrleman factor of \eqref{eq:durrleman}, evaluated on
the slice at maturity $\tau$ \citep{Gatheral2006}.
\end{proposition}

\begin{proof}
Numerator: at fixed $k$, $\partial_\tau c=\partial_w\Black\cdot
\partial_\tau w$.  With $d_\pm=-k/\sqrt w\pm\sqrt w/2$ one has
$d_+^{2}-d_-^{2}=(d_++d_-)(d_+-d_-)=-2k$, hence
$\phin(d_+)=e^{k}\phin(d_-)$, and
\[
\partial_w\Black
=\phin(d_+)\,\partial_w d_+-e^{k}\phin(d_-)\,\partial_w d_-
=\phin(d_+)\,\partial_w(d_+-d_-)
=\frac{e^{k}\phin(d_-)}{2\sqrt w}.
\]
Denominator: $\partial_{yy}c=f(\tau,y)=f_X(k)/y$ by the change of variable
$y=e^{k}$, and \eqref{eq:durrleman} gives
$f_X(k)=g_{\rm D}(k)\,\phin(d_-)/\sqrt w$.  Substituting both into
\eqref{eq:lvextract},
\[
v=\frac{e^{k}\phin(d_-)\,\partial_\tau w/\sqrt w}
       {y^{2}\,f_X(k)/y}
 =\frac{\phin(d_-)\,\partial_\tau w}{\sqrt w\,f_X(k)}
 =\frac{\partial_\tau w}{g_{\rm D}(k)}.\qedhere
\]
\end{proof}

The reader has met both halves of this ratio before, as the two families of
static validity.  The numerator is the calendar increment: nonnegative
exactly when adjacent slices are ordered in the sense of
\cref{sec:calendar}.  The denominator is positive density written in the
volatility chart, the object Chapter~3 policed slice by slice.  So
\eqref{eq:lvratio} is a compatibility theorem disguised as a formula:
\emph{a local-volatility field exists precisely where the surface is
valid}, and each mode of static arbitrage corresponds to its own sign
failure --- a calendar violation makes the numerator negative, a butterfly
violation the denominator.  This chapter therefore takes as input
hypothesis what Chapters~2--3 produced as output: slices that are
butterfly-free and calendar-ordered, in the normalized units of
\cref{sec:calendar}.  Where the hypothesis fails, no field exists and
\eqref{eq:lvratio} says so by changing sign.

\begin{figure}[htbp]
\centering
\paperfig{\textwidth}{fig_lv_ratio.pdf}
\caption{The field as a ratio, on the frozen SPY surface (stored per-expiry
fits of the snapshot, total variance interpolated linearly in $\tau$
between the eight expiries; horizontal bars mark each expiry's quoted
span).  (a)~The calendar increment $\partial_\tau w$: piecewise constant in
$\tau$ by construction, steepening toward the short-dated put side.
(b)~The Durrleman factor of each slice.  (c)~Their ratio, as local
volatility: the put wall, the quiet call side, and the short-dated
structure.  Grey cells are inadmissible (a negative numerator or
denominator): they amount to \MacLvRatioBadSharePct\% of the displayed
mesh, all in extrapolated territory --- restricted to the common quoted
support of each expiry pair the share is
\MacLvRatioBadQuotedSharePct\%, and the extracted local volatility there
ranges from \MacLvRatioLocMinPct\% to \MacLvRatioLocMaxPct\%.  Validity
where the data speaks, sign failures where conventions extrapolate: the
ratio is already a diagnostic instrument.}
\label{fig:lvratio}
\end{figure}

\Cref{fig:lvratio} evaluates the ratio on the book's frozen SPY surface ---
the eight fitted slices of the snapshot, interpolated linearly in total
variance between expiries.  The exercise makes three points at once.  The
field inherits recognizable structure from the smile: a steep put-side wall,
a quiet call side, short-dated sharpness.  The two sign conditions are
visibly the binding ones: every inadmissible cell in the picture sits
outside the common quoted support of the expiries that bracket it, in
territory owned by extrapolation conventions rather than prices.  And the
ratio was computed here from \emph{fitted, smooth} slices --- an operation
this chapter endorses.  What it must never be applied to is the market's
raw quotes, and the next section shows why.
